\scanpage{11}
is to follow. Moreover, it would be reasonable to include the fact that when $S$ is normal and the $V(a_i)$ reduced, then the $Z^{\underline{n}}_{\underline{a}}$ are normal (which now is included in your par.5), so that you can refer to this fact when mentionning examples in \deleted{sections}\inserted{par.} 2. to 5. The paragraphs 2 to 5 in this plan correspond approximately to the sections 1. to 4. in my sketchy notes. It will make some work, I am affraid, to put it into shape, but I have the feeling the notions developped with some generality can be useful still in other contexts than just for the kummer type ramification % The physical right edge follows the final n; no sentence point survives.
If you think it would be much simpler to write these sections by restricting throughout to kummer type ramification, please feel free to do so, (and refer to the possible generalization by a suitable remark). I doubt however it will be much simpler, as there are at least two useful variant\unsure{s} % Only a narrow possible s trace survives beside t at the physical right edge.
of the notion besides the one of your notes (cf ``particular case b '' in my sketch).

Some more specific observations. In theorem 6.4. (which should be formulated, I believe, in terms of tame ramification in the large sense, cf my notes), there is no need to look at the residue characteristics of the non maximal points of the $D_i$. It will turn out a posteriori (and this should be pointed out in a reamrk or corollary) that \deleted{the} if for every maximal point $s$ of every $D_i$, the inertia group at $s$ is\unsure{\tikz[baseline=0pt,x=1em,y=1ex]{\path[use as bounding box] (0,0) rectangle (.6,1.7);\fill (.17,1.03) circle[radius=.035em];\fill[rotate around={-15:(.37,1.47)}] (.37,1.47) ellipse[x radius=.035em,y radius=.07em];}} prime to the charactéeristic of $k(s)$, then the order of this group must b\unsure{\deleted{e}y|\deleted{y}e|y} automatically prime to the residue characteristic at any specialization of $s$. (This statement is non empty only when $s$ is of char. zero, of course). In fact, the proof in your notes (after reduction to the strict local case, which is immediate) works out unchanged, provided you observe that if $\unsure{AA|\lambda A|\deleted{A}A|\deleted{\lambda}A|\mathbb{A}}$ is a regular local ring, $(x_i)_{1\ i\ r}$ part of a regular system of parameters, and $(n_i)_{1\ i\ r}$ a system of intergers $\quad1$, not noecessarily
% The two small raised spots after is have unresolved origin; the inline geometry does not assert punctuation or handwriting.
% In the last letter after must, a y-descender is definite; an e overimpression and its order remain provisional.
% The faint first impression before A is present; its alphabet and any cancellation/erasure remain unresolved.
% The spaces in both lower index strings and before 1 have no visible relation signs.
% archive sheet number: 10
